\section{On-device AI}
\label{sec:on-device-ai}

All classification happens on the collar in real time, using TensorFlow Lite
Micro~\cite{david2021tflm} with the CMSIS-NN kernels~\cite{lai2018cmsisnn} for the nRF52840's
Cortex-M4F. The models are tiny int8 networks in the spirit of TinyML~\cite{warden2019tinyml}: small
enough to fit a microcontroller's RAM and run within the collar's power budget, yet accurate enough to
separate everyday behaviours.

\subsection{Confidence-gated cascade}

Two behaviours can look almost identical to a motion sensor while sounding completely different:
eating and drinking both put the head down at a bowl with similar gross motion. Rather than force one
model to separate them from accelerometry alone, Meowtion runs a confidence-gated cascade of two
models:

\begin{enumerate}[leftmargin=1.6em,itemsep=2pt]
  \item The \textbf{IMU model} runs continuously on the motion stream. It is cheap and always on, and
        it carries most of the classification load.
  \item When the IMU model is \emph{unsure} (its top class falls below a confidence threshold of
        \ph{threshold value}), the short \textbf{audio model} is invoked to confirm. Audio is therefore
        a deployed disambiguation signal, not merely a training aid.
\end{enumerate}

This keeps average power low, since the audio path runs only to break a tie, while recovering accuracy
that motion alone cannot provide. The cascade policy lives in plain C (\texttt{classifier.c}) and is
independent of the inference backend, so the gating logic can be read and tuned without touching the
TensorFlow Lite Micro code. \Cref{fig:cascade} shows the decision flow.

\begin{figure}[ht]
  \centering
  \begin{tikzpicture}[node distance=8mm and 16mm]
    \node[mwterm] (win) {IMU window};
    \node[mwflow, below=of win] (imu) {IMU model};
    \node[mwdecision, below=of imu] (q) {top class\\confident?};
    \node[mwterm, below=of q] (out1) {Report class\\+ confidence};
    \node[mwflow, right=of q] (aud) {Audio model};
    \node[mwterm, below=of aud] (out2) {Report\\confirmed class};
    \draw[mwedge] (win) -- (imu);
    \draw[mwedge] (imu) -- (q);
    \draw[mwedge] (q) -- node[left,font=\scriptsize,text=mwmut]{yes} (out1);
    \draw[mwedge] (q) -- node[above,font=\scriptsize,text=mwmut]{no} (aud);
    \draw[mwedge] (aud) -- (out2);
  \end{tikzpicture}
  \caption{The confidence-gated cascade. The IMU model runs on every window; the audio model is
           invoked only when the IMU model's top class is below the confidence threshold. With no
           model loaded the cascade returns \texttt{UNKNOWN}.}
  \label{fig:cascade}
\end{figure}

\subsection{Matching training and inference}

The on-device input transform must mirror the training transform exactly; this is the single most
important correctness property. A model trained on one representation and fed another at inference
produces confident nonsense. Meowtion enforces the match on both paths:

\begin{itemize}[leftmargin=1.5em]
  \item \textbf{Motion.} The rolling window is normalised per channel (subtract each channel's mean
        over the window, divide by the global maximum absolute value with a small epsilon) and
        quantised to int8 using the input tensor's own scale and zero-point.
  \item \textbf{Audio.} Inference audio is decimated and encoded to match the training
        representation: 8\,kHz $\mu$-law, produced by the same conversion used to prepare the training
        clips (\cref{sec:protocols}). A mismatch here was an early source of confident
        misclassification.
\end{itemize}

To avoid a large floating-point scratch buffer on a RAM-constrained part, the IMU path quantises
directly from the int16 source in a few passes (per-channel mean, then global max-abs, then write
int8) rather than materialising a float window (\cref{lst:quant}).

\begin{lstlisting}[language=C,caption={Direct int16-to-int8 quantisation, mirroring training (sketch).},label={lst:quant}]
// pass 1: per-channel mean over the window
// pass 2: global max-abs of (sample - channel_mean)
// pass 3: int8 = round( x_norm / scale ) + zero_point,
//         where x_norm = (sample - channel_mean) / (maxabs + 1e-6)
// scale and zero_point are read from the input tensor itself.
\end{lstlisting}

\subsection{Model architecture}

Each cascade stage is a small 1-D convolutional network trained in the cloud and exported as an int8
\texttt{.tflite} (\cref{sec:training-pipeline}). The IMU stage runs over a rolling window of motion
samples (six channels: accelerometer and gyroscope x/y/z); the audio stage runs over one second of
8\,kHz audio (8000 samples).

\phblock{Per-stage architecture table , for the IMU and audio networks: the input shape (IMU window
length and sample rate; audio is 8000\,$\times$\,1), the convolutional layers (filter count, kernel
size, stride), the pooling and dense layers, the output class count, and the total parameter count and
model size. This is what lets a reader reproduce or assess the networks.}

\subsection{Runtime models and model-free safety}

No model is compiled into the firmware. Each stage loads its \texttt{.tflite} at runtime
(\texttt{clf\_set\_imu\_model} / \texttt{clf\_set\_audio\_model}) from the dedicated flash partition,
delivered over the air (\cref{sec:firmware-collar,sec:protocols}), so a re-trained model swaps in
without reflashing. The engine is also safe with no model present: the presence queries return false
and the cascade returns \texttt{UNKNOWN} rather than a wrong answer. This is the state of a freshly
flashed collar before its first model arrives.

\subsection{Memory budget}

TensorFlow Lite Micro uses a fixed, statically allocated tensor arena per model (this firmware uses no
heap). The arenas dominate the RAM cost of inference and are sized to the two networks. The values in
\cref{tab:arenas} are starting budgets; each must be tuned to the real exported model, since an
undersized arena makes \texttt{AllocateTensors()} fail and the stage report itself absent.

\begin{table}[ht]
  \centering
  \caption{Tensor arenas (initial budgets).}
  \label{tab:arenas}
  \begin{tabularx}{\textwidth}{@{}llX@{}}
    \toprule
    \textbf{Model} & \textbf{Arena} & \textbf{Notes} \\
    \midrule
    IMU (slot 0)   & 16\,KiB & Always-on motion classifier, a small temporal conv net \\
    Audio (slot 1) & 48\,KiB & Larger 1-D conv over 8000 points; invoked only to confirm \\
    \bottomrule
  \end{tabularx}
\end{table}

Enabling inference initially pushed RAM use to roughly 94\,\%; eliminating the float-window scratch
buffer (quantising directly from int16) brought it back into the low 80s, leaving headroom for the BLE
stack and the raised main stack (\cref{sec:firmware-collar}).
